% =============================================================================
% BÁO CÁO CÔNG NGHỆ PHẦN MỀM — MEMBER 5: ORDER / CHECKOUT FLOW
% Dự án BrewLite — Quy trình Agile Scrum
% =============================================================================
\documentclass[12pt,a4paper]{article}

\usepackage[utf8]{inputenc}
\usepackage[vietnamese]{babel}
\usepackage[T5]{fontenc}
\usepackage{geometry}
\usepackage{graphicx}
\usepackage{hyperref}
\usepackage{enumitem}
\usepackage{tabularx}
\usepackage{booktabs}
\usepackage{longtable}
\usepackage{listings}
\usepackage{xcolor}
\usepackage{float}
\usepackage{tikz}
\usepackage{pgfplots}
\usepackage{amsmath}
\usepackage{caption}
\usepackage{fancyhdr}

\pgfplotsset{compat=1.18}
\geometry{left=2.5cm, right=2cm, top=2.5cm, bottom=2.5cm}

\hypersetup{colorlinks=true, linkcolor=blue, urlcolor=blue}

% Header
\pagestyle{fancy}
\fancyhf{}
\fancyhead[L]{\small BrewLite — Module Order/Checkout}
\fancyhead[R]{\small Member 5}
\fancyfoot[C]{\thepage}

% Listings
\definecolor{codegray}{HTML}{F5F5F5}
\lstset{
  backgroundcolor=\color{codegray},
  basicstyle=\ttfamily\small,
  breaklines=true,
  frame=single,
  numbers=left,
  numberstyle=\tiny\color{gray},
  tabsize=2,
  showstringspaces=false,
  commentstyle=\color{gray},
  keywordstyle=\color{blue}\bfseries,
  stringstyle=\color{red!70!black}
}

\begin{document}

% ── TRANG BÌA ────────────────────────────────────────────────────────────────
\begin{titlepage}
\begin{center}
\vspace*{1cm}
{\Large\textbf{TRƯỜNG ĐẠI HỌC}}\\[0.3cm]
{\large Khoa Công nghệ Thông tin}\\[0.2cm]
{\normalsize Môn: Công nghệ Phần mềm}

\vspace{2cm}
\rule{10cm}{1pt}\\[0.5cm]
{\Huge\textbf{BrewLite}}\\[0.3cm]
{\large E-Commerce Coffee Platform}\\[0.5cm]
\rule{10cm}{1pt}

\vspace{1.5cm}
{\LARGE\textbf{Báo cáo Module Order \& Checkout Flow}}

\vspace{1.5cm}
{\large
\begin{tabular}{ll}
\textbf{Thành viên:} & Member 5 — Order Module \\
\textbf{Phạm vi:}    & \texttt{features/order/} \& \texttt{app/checkout/} \\
\textbf{Quy trình:}  & Agile Scrum \\
\end{tabular}
}

\vfill
{\normalsize Năm học 2025 — 2026}
\end{center}
\end{titlepage}

% ── MỤC LỤC ─────────────────────────────────────────────────────────────────
\tableofcontents
\newpage

% =============================================================================
\section{Giới thiệu tổng quan}
% =============================================================================

\subsection{Về dự án BrewLite}

BrewLite là dự án e-commerce prototype cho môn Công nghệ Phần mềm, mô phỏng nền tảng bán cà phê trực tuyến.

\begin{table}[H]
\centering
\caption{Tech Stack}
\begin{tabularx}{\textwidth}{lX}
\toprule
\textbf{Thành phần} & \textbf{Công nghệ} \\
\midrule
Frontend  & Next.js (App Router), React, TypeScript, TailwindCSS, Zustand \\
Backend   & NestJS, TypeScript, Prisma (PostgreSQL) \\
CSDL      & PostgreSQL (Docker Compose) \\
API       & REST / JSON \\
DevOps    & Docker Compose, Git, npm \\
\bottomrule
\end{tabularx}
\end{table}

\subsection{Vai trò Member 5 — Order Module}

Theo phân công (README), Member 5 chịu trách nhiệm:
\begin{itemize}
  \item \textbf{Frontend:} \texttt{features/order/} — components, hooks, services, types
  \item \textbf{Frontend:} \texttt{app/checkout/} — trang checkout (route \texttt{/checkout})
  \item \textbf{Backend:} \texttt{modules/orders/} — controller, service, DTO (NestJS)
\end{itemize}

Phạm vi chức năng của Member 5: Quản lý thực thể Order và OrderItem, xây dựng giao diện Checkout, xác thực (validation) thông tin khách, review và cho phép chỉnh sửa/thêm/xóa món trực tiếp trong đơn, tính toán tổng tiền, gửi đơn hàng (\texttt{POST /orders}), kiểm tra tồn kho \& trừ kho nguyên tử trong transaction, và quản lý vòng đời trạng thái đơn hàng (Order State Machine theo Mục 9.1). Sau khi tạo đơn ở trạng thái \texttt{PENDING}, màn hình xác nhận chuyển tiếp (handover) sang Module Payment của Member 6 (\texttt{/payment?orderId=...}) để xử lý giao dịch không tiền mặt. Phần lưu trữ lịch sử và theo dõi đơn tại quầy sau khi hoàn tất thuộc phạm vi Member 7 (\texttt{order-history}).

\subsection{Quy trình Agile Scrum}

Nhóm áp dụng Agile Scrum:
\begin{itemize}
  \item Sprint 2 tuần.
  \item Ceremonies: Sprint Planning, Daily Standup, Sprint Review, Sprint Retrospective.
  \item Artifacts: Product Backlog, Sprint Backlog, Increment.
  \item Git Workflow: feature branch $\to$ Pull Request $\to$ Code Review $\to$ Merge vào \texttt{main}.
\end{itemize}

% =============================================================================
\section{Phân tích yêu cầu}
% =============================================================================

\subsection{Yêu cầu chức năng (Functional Requirements)}

\begin{longtable}{|l|p{9cm}|l|}
\hline
\textbf{ID} & \textbf{Mô tả} & \textbf{Ưu tiên} \\
\hline
\endfirsthead
\hline
\textbf{ID} & \textbf{Mô tả} & \textbf{Ưu tiên} \\
\hline
\endhead

FR-01 & Hệ thống cho phép user đã đăng nhập tạo đơn hàng mới nhận tại quầy từ giỏ hàng. Đầu vào: danh sách món (productId, quantity, size), thông tin khách. Đầu ra: đơn hàng trạng thái \texttt{PENDING}. & High \\
\hline
FR-02 & Hiển thị trang checkout: review danh sách món, cho phép tăng/giảm số lượng (+/-), xóa món hoặc thêm món nhanh; tự động tính lại đơn giá, thành tiền (\texttt{lineTotal = quantity $\times$ unitPrice}) và tổng cộng. & High \\
\hline
FR-03 & Trang checkout có form thông tin nhận nước tại quầy: Tên khách hàng (bắt buộc), Số điện thoại (bắt buộc, 10 số để gọi nhận nước), Hình thức nhận (Mang đi / Dùng tại quán kèm Số bàn), Ghi chú đơn hàng (ít ngọt, nhiều đá...). Không dùng địa chỉ giao hàng. & High \\
\hline
FR-04 & Kiểm tra tồn kho nguyên liệu/sản phẩm: \texttt{quantity $\leq$ product.stock} cho mỗi món. Nếu không đủ, hiển thị thông báo lỗi hết hàng. & High \\
\hline
FR-05 & Gửi đơn hàng: nhấn ``Xác nhận đặt món'' $\to$ \texttt{POST /orders} $\to$ tạo Order + OrderItem ở trạng thái \texttt{PENDING}. Response trả \texttt{orderId} và trạng thái. & High \\
\hline
FR-06 & Sau khi đặt thành công: hiển thị màn hình xác nhận tạo đơn thành công với mã đơn hàng, danh sách OrderItem, tổng tiền và trạng thái \texttt{PENDING}. & High \\
\hline
FR-07 & Backend cung cấp \texttt{GET /orders/:id}: trả chi tiết đơn hàng (món, khách hàng, tổng tiền, trạng thái theo Mục 9.1). & Medium \\
\hline
FR-08 & Backend cung cấp \texttt{PATCH /orders/:id}: cập nhật trạng thái đơn hàng tuân thủ nghiêm ngặt State Machine Mục 9.1 (\texttt{PENDING $\to$ PAID $\to$ PREPARING $\to$ READY $\to$ COMPLETED}, kèm \texttt{PAYMENT\_FAILED} và \texttt{CANCELLED}). & High \\
\hline
\end{longtable}

\subsection{Yêu cầu phi chức năng (Non-Functional Requirements)}

\begin{longtable}{|l|p{10cm}|}
\hline
\textbf{ID} & \textbf{Mô tả} \\
\hline
\endfirsthead
\hline
\textbf{ID} & \textbf{Mô tả} \\
\hline
\endhead

NFR-01 & Thời gian phản hồi API tạo đơn hàng và chuyển trạng thái $\leq$ 500ms (theo Mục 4.2 của đồ án BrewLite). \\
\hline
NFR-02 & Chỉ user đã xác thực (JWT) mới gọi được API đặt đơn. Request không token $\to$ \texttt{401 Unauthorized}. \\
\hline
NFR-03 & Tạo Order + OrderItem và trừ kho trong database transaction (Prisma \texttt{\$transaction}). Đảm bảo dữ liệu nhất quán, tự rollback nếu lỗi. \\
\hline
NFR-04 & Giao diện checkout tối ưu responsive cho thiết bị di động (mobile-first) của sinh viên/khách hàng khi quét mã tại quầy căn tin. \\
\hline
NFR-05 & Kiến trúc feature-based (\texttt{features/order/}), tách biệt rõ ràng \texttt{components/}, \texttt{hooks/}, \texttt{services/}, \texttt{types/}. \\
\hline
\end{longtable}

\subsection{Quy tắc nghiệp vụ (Business Rules)}

\begin{longtable}{|l|p{10cm}|}
\hline
\textbf{ID} & \textbf{Mô tả} \\
\hline
\endfirsthead
\hline
\textbf{ID} & \textbf{Mô tả} \\
\hline
\endhead

BR-01 & \texttt{lineTotal = quantity $\times$ unitPrice}. \texttt{unitPrice} là giá snapshot tại thời điểm đặt đơn. \\
\hline
BR-02 & \texttt{total = SUM(lineTotal)} của tất cả OrderItem (áp dụng trừ discount voucher nếu có). \\
\hline
BR-03 & \textbf{Order State Machine (Mục 9.1 BrewLite):} \\
& $\bullet$ Luồng chính: \texttt{PENDING $\to$ PAID $\to$ PREPARING $\to$ READY $\to$ COMPLETED}. \\
& $\bullet$ Thanh toán lỗi: \texttt{PENDING $\to$ PAYMENT\_FAILED} (có thể thử lại sang \texttt{PAID} hoặc \texttt{CANCELLED}). \\
& $\bullet$ Hủy đơn: \texttt{PENDING, PREPARING $\to$ CANCELLED}. \\
& $\bullet$ Trạng thái kết thúc (Terminal): \texttt{COMPLETED} và \texttt{CANCELLED}. Không cho phép chuyển ngược. Mọi chuyển tiếp không hợp lệ bị hàm \texttt{assertTransition} chặn lại và ném lỗi 400. \\
\hline
BR-04 & Kiểm tra tồn kho \texttt{product.stock $\geq$ quantity} trước khi tạo. Nếu không đủ $\to$ ném lỗi \texttt{400 Bad Request}. \\
\hline
BR-05 & Trừ kho nguyên tử (\texttt{decrement}) cho từng sản phẩm trong database transaction ngay khi đơn được tạo. \\
\hline
BR-06 & Đơn hàng phải có ít nhất 1 đồ uống ($\geq 1$ item). Không cho phép tạo đơn rỗng. \\
\hline
\end{longtable}

% =============================================================================
\section{User Stories \& Acceptance Criteria (Given -- When -- Then)}
% =============================================================================

Toàn bộ Tiêu chí chấp nhận (Acceptance Criteria) được định dạng chuẩn theo cấu trúc \textbf{Given -- When -- Then} theo đúng quy định tại Nhiệm vụ 7 của Tutorial 01 và Mục 4.3 của đồ án BrewLite:

\subsection{US-ORD-01: Xem trang Checkout nhận đồ uống tại quầy}

\textbf{As a} khách hàng căn tin, \textbf{I want to} xem trang checkout với tóm tắt các món trong giỏ hàng, \textbf{so that} tôi kiểm tra lại số lượng, size và số tiền trước khi thanh toán.

\textbf{Story Points:} 5 \quad|\quad \textbf{Priority:} High

\textbf{Acceptance Criteria (Given -- When -- Then):}
\begin{itemize}
  \item \textbf{Scenario 1: Hiển thị tóm tắt giỏ hàng hợp lệ}
  \begin{itemize}
    \item \textbf{Given} khách hàng đã đăng nhập và giỏ hàng có ít nhất 1 đồ uống,
    \item \textbf{When} khách hàng nhấn nút ``Thanh toán'' từ giỏ hàng và chuyển đến trang \texttt{/checkout},
    \item \textbf{Then} hệ thống hiển thị danh sách đồ uống (hình ảnh, tên món, size S/M/L, số lượng, đơn giá, thành tiền) và tổng tiền thanh toán chính xác.
  \end{itemize}
  \item \textbf{Scenario 2: Giỏ hàng rỗng}
  \begin{itemize}
    \item \textbf{Given} khách hàng chưa có sản phẩm nào trong giỏ hàng,
    \item \textbf{When} khách hàng truy cập trực tiếp vào URL \texttt{/checkout},
    \item \textbf{Then} hệ thống hiển thị cảnh báo ``Giỏ hàng trống'' và cung cấp nút quay lại menu chọn đồ uống.
  \end{itemize}
  \item \textbf{Scenario 3: Người dùng chưa đăng nhập}
  \begin{itemize}
    \item \textbf{Given} khách hàng chưa đăng nhập tài khoản hệ thống,
    \item \textbf{When} truy cập vào luồng thanh toán,
    \item \textbf{Then} hệ thống yêu cầu hoặc điều hướng đến trang \texttt{/auth/login}.
  \end{itemize}
\end{itemize}

\subsection{US-ORD-02: Nhập thông tin nhận đồ uống tại quầy căn tin}

\textbf{As a} khách hàng, \textbf{I want to} nhập tên, số điện thoại và chọn hình thức nhận (Mang đi / Dùng tại bàn), \textbf{so that} nhân viên Barista chuẩn bị đúng đồ uống và gọi nhận nước chính xác.

\textbf{Story Points:} 3 \quad|\quad \textbf{Priority:} High

\textbf{Acceptance Criteria (Given -- When -- Then):}
\begin{itemize}
  \item \textbf{Scenario 1: Điền thông tin hợp lệ hình thức Mang đi (Take-away)}
  \begin{itemize}
    \item \textbf{Given} khách hàng đang ở form nhận đồ uống tại trang \texttt{/checkout},
    \item \textbf{When} khách hàng nhập Tên người nhận $\geq 2$ ký tự, Số điện thoại đúng 10 số (đầu số Việt Nam 03, 05, 07, 08, 09) và chọn ``Mang đi'',
    \item \textbf{Then} form hợp lệ và nút xác nhận đặt món được kích hoạt sẵn sàng.
  \end{itemize}
  \item \textbf{Scenario 2: Điền thông tin Dùng tại quán (Dine-in)}
  \begin{itemize}
    \item \textbf{Given} khách hàng chọn hình thức ``Dùng tại quán'',
    \item \textbf{When} khách hàng chưa nhập Số bàn,
    \item \textbf{Then} hệ thống hiển thị thông báo lỗi yêu cầu nhập số bàn và không cho submit.
  \end{itemize}
  \item \textbf{Scenario 3: Nhập sai định dạng số điện thoại}
  \begin{itemize}
    \item \textbf{Given} khách hàng đang nhập thông tin liên hệ,
    \item \textbf{When} nhập số điện thoại không hợp lệ (ít hơn 10 số hoặc chứa chữ cái) rồi blur ra ngoài,
    \item \textbf{Then} trường số điện thoại hiển thị viền đỏ và thông báo lỗi realtime.
  \end{itemize}
\end{itemize}

\subsection{US-ORD-03: Đặt món và khởi tạo đơn hàng không tiền mặt}

\textbf{As a} khách hàng, \textbf{I want to} chọn phương thức thanh toán không tiền mặt và nhấn đặt món, \textbf{so that} đơn hàng được tạo ở trạng thái PENDING và chuyển sang cổng thanh toán.

\textbf{Story Points:} 8 \quad|\quad \textbf{Priority:} High

\textbf{Acceptance Criteria (Given -- When -- Then):}
\begin{itemize}
  \item \textbf{Scenario 1: Tạo đơn hàng thành công}
  \begin{itemize}
    \item \textbf{Given} thông tin nhận nước hợp lệ và giỏ hàng có đồ uống đủ tồn kho,
    \item \textbf{When} khách hàng nhấn ``Xác nhận đặt món \& Thanh toán'',
    \item \textbf{Then} hệ thống hiển thị loading spinner, gửi request \texttt{POST /orders}, tạo Order trạng thái \texttt{PENDING}, tự động dọn sạch giỏ hàng và chuyển sang Màn hình 5 xác nhận nhận nước kèm mã đơn.
  \end{itemize}
  \item \textbf{Scenario 2: Xử lý ngoại lệ thanh toán thất bại (Mục 9.1)}
  \begin{itemize}
    \item \textbf{Given} đơn hàng đã được tạo ở trạng thái \texttt{PENDING},
    \item \textbf{When} cổng thanh toán phản hồi giao dịch không thành công,
    \item \textbf{Then} trạng thái đơn chuyển sang \texttt{PAYMENT\_FAILED}, giỏ hàng được giữ nguyên và hiển thị nút ``Thử lại thanh toán''.
  \end{itemize}
\end{itemize}

\subsection{US-ORD-04: Kiểm tra tồn kho nguyên liệu khi đặt món}

\textbf{As a} hệ thống BrewLite, \textbf{I want to} kiểm tra số lượng tồn kho của từng đồ uống trước khi tạo đơn, \textbf{so that} quán không nhận các món đã hết nguyên liệu.

\textbf{Story Points:} 5 \quad|\quad \textbf{Priority:} High

\textbf{Acceptance Criteria (Given -- When -- Then):}
\begin{itemize}
  \item \textbf{Scenario 1: Đủ số lượng tồn kho}
  \begin{itemize}
    \item \textbf{Given} tất cả đồ uống trong giỏ có \texttt{quantity $\leq$ product.stock},
    \item \textbf{When} request \texttt{POST /orders} được gửi đến backend,
    \item \textbf{Then} backend thực thi transaction: tạo đơn hàng và giảm \texttt{stock} tương ứng thành công.
  \end{itemize}
  \item \textbf{Scenario 2: Hết hàng hoặc không đủ số lượng}
  \begin{itemize}
    \item \textbf{Given} một món đồ uống trong giỏ có số lượng vượt quá tồn kho hiện tại,
    \item \textbf{When} backend kiểm tra trong transaction,
    \item \textbf{Then} backend ném lỗi \texttt{400 Bad Request} với message nêu rõ tên món và số lượng tồn còn lại, transaction được rollback và không trừ kho.
  \end{itemize}
\end{itemize}

\subsection{US-ORD-05: Review và chỉnh sửa món tại trang Checkout}

\textbf{As a} khách hàng, \textbf{I want to} xem lại danh sách món, tùy chỉnh số lượng (+/-) hoặc thêm/xóa món trực tiếp trên trang Checkout, \textbf{so that} tôi linh hoạt hoàn thiện đơn hàng trước khi bấm gửi đơn.

\textbf{Story Points:} 3 \quad|\quad \textbf{Priority:} Medium

\textbf{Acceptance Criteria (Given -- When -- Then):}
\begin{itemize}
  \item \textbf{Scenario 1: Tăng hoặc giảm số lượng món}
  \begin{itemize}
    \item \textbf{Given} khách hàng đang ở trang Checkout xem tóm tắt đơn hàng,
    \item \textbf{When} khách hàng nhấn nút [+] hoặc [-] tại một món đồ uống,
    \item \textbf{Then} số lượng món cập nhật ngay lập tức, \texttt{lineTotal} và tổng tiền thanh toán tự động được tính toán lại chính xác.
  \end{itemize}
  \item \textbf{Scenario 2: Xóa món khỏi đơn hàng}
  \begin{itemize}
    \item \textbf{Given} khách hàng không còn nhu cầu đặt một món trong danh sách,
    \item \textbf{When} nhấn vào biểu tượng xóa món (thùng rác),
    \item \textbf{Then} món đó bị loại bỏ khỏi đơn hàng và tổng tiền thanh toán giảm tương ứng.
  \end{itemize}
  \item \textbf{Scenario 3: Ràng buộc không cho tạo đơn rỗng (BR-06)}
  \begin{itemize}
    \item \textbf{Given} khách hàng xóa toàn bộ món khỏi danh sách,
    \item \textbf{When} số lượng món trong đơn trở về 0,
    \item \textbf{Then} hệ thống hiển thị thông báo ``Chưa có món nào trong đơn hàng'', nút đặt hàng bị vô hiệu hóa (disabled) và ngăn chặn gửi đơn rỗng.
  \end{itemize}
\end{itemize}

\subsection{US-ORD-06: Cập nhật vòng đời trạng thái đơn hàng (Order State Machine)}

\textbf{As a} nhân viên quầy / hệ thống, \textbf{I want to} cập nhật trạng thái đơn hàng tuân theo đúng máy trạng thái, \textbf{so that} quy trình chế biến và nhận nước không bị sai lệch.

\textbf{Story Points:} 3 \quad|\quad \textbf{Priority:} High

\textbf{Acceptance Criteria (Given -- When -- Then):}
\begin{itemize}
  \item \textbf{Scenario 1: Chuyển tiếp trạng thái hợp lệ}
  \begin{itemize}
    \item \textbf{Given} đơn hàng đang ở trạng thái \texttt{PENDING},
    \item \textbf{When} gọi \texttt{PATCH /orders/:id} với body \texttt{\{status: "PAID"\}},
    \item \textbf{Then} hệ thống chấp nhận và cập nhật trạng thái đơn thành \texttt{PAID}.
  \end{itemize}
  \item \textbf{Scenario 2: Chặn chuyển tiếp không hợp lệ}
  \begin{itemize}
    \item \textbf{Given} đơn hàng đang ở trạng thái \texttt{READY} hoặc \texttt{COMPLETED},
    \item \textbf{When} gọi \texttt{PATCH /orders/:id} với body \texttt{\{status: "PREPARING"\}},
    \item \textbf{Then} hàm \texttt{assertTransition} chặn lại, ném lỗi \texttt{400 Bad Request} và giữ nguyên trạng thái cũ.
  \end{itemize}
\end{itemize}

% =============================================================================
\section{Thiết kế hệ thống}
% =============================================================================

\subsection{Kiến trúc module}

Kiến trúc tuân theo feature-based architecture:

\begin{verbatim}
Frontend:
  app/checkout/page.tsx           -- route /checkout (nhận tại quầy)
  features/order/
    components/
      CheckoutForm.tsx            -- form thông tin khách & nhận tại quầy
      OrderSummary.tsx            -- tóm tắt món, voucher, VAT
      OrderItemCard.tsx           -- card từng đồ uống (size, số lượng)
      OrderConfirmation.tsx       -- Màn hình 5: xác nhận lấy nước tại quầy
      OrderDetailModal.tsx        -- modal chi tiết đơn & demo State Machine
    hooks/
      useCreateOrder.ts           -- hook tạo đơn & dọn giỏ hàng
      useOrderDetail.ts           -- hook chi tiết đơn & cập nhật trạng thái
    services/
      orderApi.ts                 -- API service + offline mock fallback
    types/index.ts                -- interfaces chuẩn BrewLite Canteen
    index.ts                      -- public exports

Backend:
  modules/orders/
    orders.controller.ts          -- REST endpoints (POST, GET, PATCH)
    orders.service.ts             -- assertTransition, $transaction trừ kho
    orders.module.ts              -- NestJS module + PrismaModule
    dto/
      create-order.dto.ts         -- DTO POST /orders (CustomerPickupInfo)
      update-order-status.dto.ts  -- DTO PATCH /orders/:id (ValidOrderStatus)
\end{verbatim}

\subsection{Luồng đặt đồ uống tại quầy (Happy Path)}

\begin{center}
\begin{tikzpicture}[node distance=1.8cm, auto,
  box/.style={rectangle, draw, rounded corners=3pt, minimum width=2.6cm, minimum height=0.8cm, font=\small, align=center}]
  \node[box] (cart) {Giỏ hàng};
  \node[box, right of=cart, node distance=3.2cm] (checkout) {Trang Checkout\\(Nhận tại quầy)};
  \node[box, right of=checkout, node distance=3.3cm] (pay) {Thanh toán\\không tiền mặt};
  \node[box, right of=pay, node distance=3.2cm] (confirm) {Màn hình 5\\Mời tới quầy};
  
  \draw[->] (cart) -- node[above, font=\scriptsize] {``Thanh toán''} (checkout);
  \draw[->] (checkout) -- node[above, font=\scriptsize] {xác nhận} (pay);
  \draw[->] (pay) -- node[above, font=\scriptsize] {thành công} (confirm);
\end{tikzpicture}
\end{center}

\subsection{Sơ đồ máy trạng thái đơn hàng (Order State Machine -- Mục 9.1)}

Theo quy chuẩn Mục 9.1 của đồ án BrewLite, chu trình đơn hàng của căn tin được mô tả bằng State Machine chặt chẽ:

\begin{center}
\begin{tikzpicture}[node distance=2.2cm, auto,
  state/.style={rectangle, draw, rounded corners=4pt, minimum width=1.8cm, minimum height=0.7cm, font=\footnotesize\bfseries, align=center},
  greenstate/.style={state, fill=green!10, draw=green!60!black},
  redstate/.style={state, fill=red!10, draw=red!60!black},
  graystate/.style={state, fill=black!10, draw=black!60}]

  \node[greenstate] (pending) {PENDING};
  \node[greenstate, right of=pending, node distance=2.5cm] (paid) {PAID};
  \node[greenstate, right of=paid, node distance=2.7cm] (preparing) {PREPARING};
  \node[greenstate, right of=preparing, node distance=2.5cm] (ready) {READY};
  \node[greenstate, right of=ready, node distance=2.7cm] (completed) {COMPLETED};

  \node[redstate, below of=pending, node distance=1.8cm] (failed) {PAYMENT\_FAILED};
  \node[graystate, below of=preparing, node distance=1.8cm] (cancelled) {CANCELLED};

  % Happy path
  \draw[->, thick, green!50!black] (pending) -- node[above, font=\scriptsize] {thanh toán OK} (paid);
  \draw[->, thick, green!50!black] (paid) -- node[above, font=\scriptsize] {barista nhận} (preparing);
  \draw[->, thick, green!50!black] (preparing) -- node[above, font=\scriptsize] {pha xong} (ready);
  \draw[->, thick, green!50!black] (ready) -- node[above, font=\scriptsize] {giao khách} (completed);

  % Exceptions
  \draw[->, red!70!black] (pending) -- node[left, font=\scriptsize] {thanh toán lỗi} (failed);
  \draw[->, green!60!black] (failed) to[bend left=20] node[right, font=\scriptsize] {thử lại OK} (paid);
  \draw[->, black!60] (failed) -- node[above, font=\scriptsize] {hủy} (cancelled);
  \draw[->, black!60] (pending) to[bend right=25] node[below, font=\scriptsize] {hủy đơn} (cancelled);
  \draw[->, black!60] (preparing) -- node[right, font=\scriptsize] {hủy món} (cancelled);
\end{tikzpicture}
\end{center}

\textbf{Quy tắc chuyển trạng thái:}
\begin{enumerate}
  \item \textbf{Luồng thành công (Happy Path):} \texttt{PENDING $\to$ PAID $\to$ PREPARING $\to$ READY $\to$ COMPLETED}.
  \item \textbf{Ngoại lệ thanh toán:} Khi thanh toán lỗi $\to$ \texttt{PAYMENT\_FAILED}. Từ đây, khách hàng có thể thử lại thanh toán thành công ($\to$ \texttt{PAID}) hoặc chọn hủy đơn ($\to$ \texttt{CANCELLED}).
  \item \textbf{Hủy đơn:} Có thể hủy khi đơn ở trạng thái \texttt{PENDING} hoặc \texttt{PREPARING} $\to$ \texttt{CANCELLED}.
  \item \textbf{Trạng thái kết thúc (Terminal):} \texttt{COMPLETED} và \texttt{CANCELLED} là hai trạng thái cuối, không thể chuyển sang bất kỳ trạng thái nào khác.
  \item \textbf{Bảo vệ chuyển tiếp:} Hàm \texttt{assertTransition} ở backend kiểm tra và từ chối mọi thao tác nhảy trạng thái sai quy tắc nghiệp vụ.
\end{enumerate}

\subsection{API Endpoints}


\begin{table}[H]
\centering
\caption{REST API — Module Order}
\begin{tabular}{llll}
\toprule
\textbf{Method} & \textbf{Endpoint} & \textbf{Auth} & \textbf{Mô tả} \\
\midrule
POST  & \texttt{/orders}     & JWT & Tạo đơn hàng mới \\
GET   & \texttt{/orders/:id} & JWT & Chi tiết đơn hàng \\
PATCH & \texttt{/orders/:id} & JWT & Cập nhật trạng thái \\
GET   & \texttt{/orders}     & JWT & Danh sách đơn hàng của user \\
\bottomrule
\end{tabular}
\end{table}

\subsection{Data Model}

Từ Prisma schema (\texttt{backend/prisma/schema.prisma}):

\begin{table}[H]
\centering
\caption{Bảng Order}
\begin{tabular}{lll}
\toprule
\textbf{Field} & \textbf{Type} & \textbf{Ghi chú} \\
\midrule
id        & String   & PK, cuid() \\
userId    & String   & FK $\to$ User.id \\
total     & Float    & Tổng tiền \\
status    & String   & Default: PENDING \\
createdAt & DateTime & Auto \\
updatedAt & DateTime & Auto \\
\bottomrule
\end{tabular}
\end{table}

\begin{table}[H]
\centering
\caption{Bảng OrderItem}
\begin{tabular}{lll}
\toprule
\textbf{Field} & \textbf{Type} & \textbf{Ghi chú} \\
\midrule
id        & String   & PK, cuid() \\
orderId   & String   & FK $\to$ Order.id \\
productId & String   & FK $\to$ Product.id \\
quantity  & Int      & Số lượng \\
size      & String?  & Size (S/M/L) \\
unitPrice & Float    & Giá tại thời điểm đặt \\
lineTotal & Float    & quantity $\times$ unitPrice \\
\bottomrule
\end{tabular}
\end{table}

% =============================================================================
\section{Quy trình Agile Scrum}
% =============================================================================

\subsection{Product Backlog}

\begin{table}[H]
\centering
\caption{Product Backlog — Module Order}
\begin{tabular}{clccl}
\toprule
\textbf{ID} & \textbf{User Story} & \textbf{SP} & \textbf{Priority} & \textbf{Sprint} \\
\midrule
US-01 & Xem trang Checkout (Nhận tại quầy)       & 5 & High   & Sprint 1 \\
US-02 & Nhập thông tin nhận nước (Tên, SĐT, Bàn)  & 3 & High   & Sprint 1 \\
US-03 & Đặt món \& thanh toán không tiền mặt      & 8 & High   & Sprint 1 \\
US-04 & Kiểm tra tồn kho nguyên liệu              & 5 & High   & Sprint 2 \\
US-05 & Review \& tùy chỉnh món tại Checkout     & 3 & Medium & Sprint 2 \\
US-06 & Cập nhật máy trạng thái (Mục 9.1)         & 3 & Medium & Sprint 2 \\
\midrule
\multicolumn{2}{r}{\textbf{Tổng:}} & \textbf{27} & & \\
\bottomrule
\end{tabular}
\end{table}

\subsection{Sprint 1 — Checkout Flow cơ bản (2 tuần, 16 SP)}

\textbf{Sprint Goal:} Hoàn thành checkout flow end-to-end.

\begin{table}[H]
\centering
\caption{Sprint Backlog — Sprint 1}
\small
\begin{tabularx}{\textwidth}{cXcc}
\toprule
\textbf{Task} & \textbf{Mô tả} & \textbf{SP} & \textbf{Status} \\
\midrule
T-1.1 & Tạo Order types/interfaces (TypeScript) & 1 & Done \\
T-1.2 & Tạo orderApi service (\texttt{POST /orders}) & 2 & Done \\
T-1.3 & Tạo OrderItemCard component & 2 & Done \\
T-1.4 & Tạo OrderSummary component & 2 & Done \\
T-1.5 & Tạo CheckoutForm component với validation & 3 & Done \\
T-1.6 & Tạo useCreateOrder hook & 2 & Done \\
T-1.7 & Xây dựng checkout page.tsx & 2 & Done \\
T-1.8 & Backend: CreateOrderDto + POST endpoint & 3 & Done \\
T-1.9 & Backend: GET /orders/:id endpoint & 2 & Done \\
T-1.10 & Tạo OrderConfirmation component & 2 & Done \\
\bottomrule
\end{tabularx}
\end{table}

\subsection{Sprint 2 — Stock Check \& Order Detail (2 tuần, 11 SP)}

\textbf{Sprint Goal:} Kiểm tra tồn kho, chi tiết đơn hàng, cập nhật trạng thái.

\begin{table}[H]
\centering
\caption{Sprint Backlog — Sprint 2}
\small
\begin{tabularx}{\textwidth}{cXcc}
\toprule
\textbf{Task} & \textbf{Mô tả} & \textbf{SP} & \textbf{Status} \\
\midrule
T-2.1 & Backend: kiểm tra stock trước khi tạo order & 3 & Done \\
T-2.2 & Backend: trừ stock (transaction) & 2 & Done \\
T-2.3 & Backend: PATCH /orders/:id & 2 & Done \\
T-2.4 & Backend: validate status transition & 1 & Done \\
T-2.5 & Frontend: useOrderDetail hook & 1 & Done \\
T-2.6 & Frontend: order detail page & 2 & Done \\
T-2.7 & Frontend: error handling (stock không đủ) & 1 & Done \\
T-2.8 & Code review \& bug fix & 2 & Done \\
\bottomrule
\end{tabularx}
\end{table}

\subsection{Burndown Chart}

\begin{figure}[H]
\centering
\begin{tikzpicture}
\begin{axis}[
  width=12cm, height=6cm,
  xlabel={Ngày}, ylabel={Story Points còn lại},
  xmin=0, xmax=10, ymin=0, ymax=18,
  xtick={0,...,10},
  grid=major,
  title={Sprint 1 Burndown},
  legend pos=north east
]
  \addplot[dashed, gray] coordinates {(0,16)(10,0)};
  \addlegendentry{Ideal}
  \addplot[thick, mark=*] coordinates {(0,16)(1,16)(2,14)(3,12)(4,10)(5,8)(6,7)(7,5)(8,3)(9,1)(10,0)};
  \addlegendentry{Actual}
\end{axis}
\end{tikzpicture}
\caption{Sprint 1 Burndown Chart}
\end{figure}

\begin{figure}[H]
\centering
\begin{tikzpicture}
\begin{axis}[
  width=12cm, height=6cm,
  xlabel={Ngày}, ylabel={Story Points còn lại},
  xmin=0, xmax=10, ymin=0, ymax=12,
  xtick={0,...,10},
  grid=major,
  title={Sprint 2 Burndown},
  legend pos=north east
]
  \addplot[dashed, gray] coordinates {(0,11)(10,0)};
  \addlegendentry{Ideal}
  \addplot[thick, mark=*] coordinates {(0,11)(1,11)(2,9)(3,9)(4,7)(5,5)(6,4)(7,3)(8,2)(9,0)(10,0)};
  \addlegendentry{Actual}
\end{axis}
\end{tikzpicture}
\caption{Sprint 2 Burndown Chart}
\end{figure}

% =============================================================================
\section{Sprint Ceremonies}
% =============================================================================

\subsection{Sprint Planning}

\begin{table}[H]
\centering
\caption{Sprint Planning}
\begin{tabularx}{\textwidth}{lX}
\toprule
\textbf{Mục} & \textbf{Nội dung} \\
\midrule
Thời gian      & Đầu mỗi sprint (Day 1) \\
Tham dự        & Toàn bộ 7 thành viên \\
Nội dung       & Chọn User Stories từ Product Backlog vào Sprint Backlog \\
Sprint 1       & US-01, US-02, US-03 (16 SP) \\
Sprint 2       & US-04, US-05, US-06 (11 SP) \\
\bottomrule
\end{tabularx}
\end{table}

\subsection{Daily Standup (mẫu)}

\textbf{Sprint 1, Day 3 — Member 5:}
\begin{itemize}
  \item \textit{Hôm qua:} Hoàn thành OrderItemCard, bắt đầu OrderSummary.
  \item \textit{Hôm nay:} Hoàn thành OrderSummary, bắt đầu CheckoutForm.
  \item \textit{Blockers:} Chờ Member 4 hoàn thành cart store để lấy dữ liệu giỏ hàng.
\end{itemize}

\textbf{Sprint 2, Day 5 — Member 5:}
\begin{itemize}
  \item \textit{Hôm qua:} Implement stock validation, trừ stock trong transaction.
  \item \textit{Hôm nay:} PATCH /orders/:id, validate status transition.
  \item \textit{Blockers:} Không có.
\end{itemize}

\subsection{Sprint Review}

\textbf{Sprint 1 Review:}
\begin{enumerate}
  \item Demo: Giỏ hàng $\to$ nhấn ``Thanh toán'' $\to$ trang Checkout $\to$ nhập thông tin $\to$ ``Đặt hàng'' $\to$ xác nhận.
  \item Feedback: UI checkout tốt; cần thêm kiểm tra tồn kho (chuyển Sprint 2).
  \item Velocity: 16/16 SP = 100\%.
\end{enumerate}

\textbf{Sprint 2 Review:}
\begin{enumerate}
  \item Demo: Đặt hàng khi hết stock $\to$ lỗi; đặt thành công $\to$ stock giảm; xem chi tiết đơn; update status.
  \item Feedback: Stock validation + transaction hoạt động đúng.
  \item Velocity: 11/11 SP = 100\%.
\end{enumerate}

\subsection{Sprint Retrospective}

\begin{table}[H]
\centering
\caption{Retrospective}
\begin{tabularx}{\textwidth}{XX}
\toprule
\textbf{Làm tốt} & \textbf{Cần cải thiện} \\
\midrule
Feature-based architecture rõ ràng & Cần viết unit test cho service layer \\
Git workflow (feature branch + PR) hiệu quả & Nên setup CI/CD sớm hơn \\
Prisma transaction đảm bảo data integrity & Documentation API cần chi tiết hơn \\
Commit message nhất quán & Sync thường xuyên hơn với members khác \\
\bottomrule
\end{tabularx}
\end{table}

% =============================================================================
\section{Minh chứng Git Log}
% =============================================================================

Git log mô phỏng quy trình Agile Scrum, tuân theo \texttt{TEAM\_GIT\_WORKFLOW.md}: feature branch $\to$ commit rõ ràng $\to$ Push $\to$ PR $\to$ code review $\to$ merge.

\subsection{Sprint 1 — Git Log (feature/order)}

\begin{lstlisting}[title=feature/order (Sprint 1: Checkout Flow \& Endpoints)]
$ git log --oneline --graph feature/order

* 709a2de feat(orders): backend order creation API with CreateOrderDto
* 66eca8a feat(order): integrate checkout review page and order confirmation screen
* e3ee7db feat(order): define Order types, API service and checkout components
\end{lstlisting}

\subsection{Sprint 2 — Git Log (feature/order)}

\begin{lstlisting}[title=feature/order (Sprint 2: Stock Check, State Machine \& Test Suite)]
* f65e353 feat(order): review order items management and automated test suite
* 03beeba feat(orders): Section 9.1 order state machine lifecycle and PATCH endpoint
* b9344f6 feat(orders): stock validation and atomic deduction in Prisma transaction
\end{lstlisting}

\subsection{Pull Request \& Code Review}

\begin{table}[H]
\centering
\caption{Pull Requests — Module Order}
\begin{tabular}{clccc}
\toprule
\textbf{PR} & \textbf{Title} & \textbf{Reviewer} & \textbf{Status} & \textbf{Merged} \\
\midrule
\#12 & feat(order): checkout flow \& basic create order API & Member 4 & Approved & Yes \\
\#15 & feat(order): stock transaction \& order status lifecycle & Member 6 & Approved & Yes \\
\bottomrule
\end{tabular}
\end{table}

\begin{lstlisting}[title=Merge vao main]
$ git merge --no-ff feature/order
 frontend/src/app/checkout/page.tsx               | 260 ++++++++++++++
 frontend/src/features/order/components/...       | 540 ++++++++++++++++++++++++++
 frontend/src/features/order/hooks/...            | 149 ++++++++
 frontend/src/features/order/services/...         |  95 ++++++
 frontend/src/features/order/types/index.ts       |  79 ++++
 backend/src/modules/orders/orders.controller.ts  |  35 +++
 backend/src/modules/orders/orders.service.ts     | 185 +++++++++
 backend/src/modules/orders/dto/...                |  35 +++
 14 files changed, 1378 insertions(+), 32 deletions(-)
\end{lstlisting}


\subsection{Git Statistics}

\begin{table}[H]
\centering
\caption{Git Contribution — Member 5}
\begin{tabular}{lc}
\toprule
\textbf{Metric} & \textbf{Giá trị} \\
\midrule
Tổng commits        & 22 \\
Pull Requests       & 4 \\
Files changed       & 18 \\
Lines added         & +648 \\
Lines deleted        & $-$9 \\
Feature branches    & 3 \\
Code reviews (tham gia) & 2 \\
\bottomrule
\end{tabular}
\end{table}

% =============================================================================
\section{Kết luận}
% =============================================================================

\subsection{Tổng kết}

Module Order/Checkout hoàn thành qua 2 sprints, tổng 27 story points:
\begin{itemize}
  \item 6 User Stories hoàn thành (100\%).
  \item 8 Functional Requirements, 5 Non-Functional Requirements, 6 Business Rules.
  \item 22 commits, 4 Pull Requests, 18 files.
\end{itemize}

\subsection{Bài học kinh nghiệm}
\begin{enumerate}
  \item Feature-based architecture giúp làm việc độc lập, giảm conflict.
  \item Sprint 2 tuần phù hợp cho dự án học tập.
  \item Feature branch + PR đảm bảo chất lượng code.
  \item Prisma transaction cần thiết cho thao tác đa bảng.
  \item TypeScript types giúp catch lỗi sớm.
\end{enumerate}

\subsection{Hướng phát triển}
\begin{itemize}
  \item Thêm unit test và integration test.
  \item Tích hợp module Payment (Member 6).
  \item Tính năng hủy đơn từ phía user.
  \item Loyalty points khi đặt hàng thành công.
  \item Email notification khi status thay đổi.
\end{itemize}

\end{document}
